\documentclass{article}
\usepackage{amsmath,amssymb}
\begin{document}
\paragraph{Registered hybrid carrier dynamics and subthreshold state-collision fixture.}
Protocol~V2.6 defines the complete carrier state as
$\mathbf{z}_n=(u_n,A_n,r_n)$, with $r_n$ the remaining refractory time,
not a surprise-weighted accumulator. The registered implementation uses
Euler steps $h=0.1\,\text{ms}$, $\tau_m=10\,\text{ms}$,
$\tau_A=15\,\text{ms}$, $\tau_{\mathrm{ref}}=2\,\text{ms}$,
threshold $1$, and rest/reset potential $0$. For M1,
$\mathbf{z}_0=(0,0,0)$, and
\begin{equation}
\begin{aligned}
 A_{n+1}&=\min\!\left(5,\max\!\left(0,
 A_n+\frac{h}{\tau_A}(-A_n+\nu_n)\right)\right),\\
 u_n^\star&=u_n+\frac{h}{\tau_m}
 \left(-u_n+A_n I_{\mathrm{content},n}+I_{\mathrm{query},n}\right)
 \quad (r_n=0).
\end{aligned}
\label{eq:p1b_carrier}
\end{equation}
The membrane update uses the \emph{old} gate $A_n$; the gate evolves even
during refractoriness. The remaining updates are
\[
(u_{n+1},r_{n+1},s_{n+1})=
\begin{cases}
 (0,\max(0,r_n-h),0),&r_n>0,\\
 (0,\tau_{\mathrm{ref}},1),&r_n=0,\ u_n^\star\ge1,\\
 (u_n^\star,0,0),&r_n=0,\ u_n^\star<1.
\end{cases}
\]
For an active event $(K_i,V_i)$, the registered unit weights give
$I_{\mathrm{content}}=K_i+V_i$ and
$\nu=\max(0,|K_i|-1.4)$; both drives are zero during blanks, ISIs and
query presentation. A query supplies $I_{\mathrm{query}}=q$ for
$10\,\text{ms}$, and is zero otherwise. The two roles independently probe
the same encoded and delayed checkpoint.

For the analog decoder, use cue-local indices $0,\ldots,100$.
At the deadline, $u_{\mathrm{readout}}=0$ if $r_{99}>0$; otherwise
$u_{\mathrm{readout}}=u_{99}^\star$, including when that last step itself
crosses threshold and resets the stored voltage. The frozen readout is
$\hat v=g_{\mathrm{dec}}u_{\mathrm{readout}}+b_{\mathrm{dec}}$,
with displayed rounded coefficients $g_{\mathrm{dec}}=0.132630817$ and
$b_{\mathrm{dec}}=-0.026816484$. The query affects the preceding dynamics
but is not a separate argument of this affine decoder. The coefficients
were calibrated once on 500 independent histories and shared across
the registered analog arms.

Fixture~G1.5 is a specific compensated-value construction on a verified
subthreshold branch, not a claim that all subthreshold inputs are erased.
It fixes positive-novelty keys $K_1=1.8$, $K_2=2.2$ in slots 0 and 1
and uses the fixture-specific content weights $w_K=w_V=0.1$ rather than
the main experiment's unit weights. Since $A_n$ depends on the fixed keys
and timing and no spikes occur, $u_T=b+c_1V_1+c_2V_2$ and $r_T=0$.
Choosing $\Delta V_1=2.5$ and
$\Delta V_2=-(c_1/c_2)\Delta V_1$ gives equal ideal terminal states;
the executable checks legal values, no spikes and a floating residual
$<10^{-10}$. The common query $q=1.8$ requests different first values,
giving a pairwise worst-case ideal-state delivery error of at least
$1.25>0.50$. This is separate from Protocol~V2.6's first-event
zero-novelty erasure result under $A(0)=0$.

\paragraph{Confirmatory evaluation: floor-limited comparison and unestablished redundancy.}
The confirmatory audit spanned 4,000 unique histories across 3 post-sequence silence
delays ($\Delta t_{\text{silence}}\in\{0,5,10\}\,\text{ms}$) and 8 model variants,
totaling 96,000 model-episodes (192,000 role queries). Operand delivery requires
$|\hat V-V_{\mathrm{target}}|<0.50$ without abstention; joint delivery requires
both roles to succeed. Protocol~V2.6 \S7.3 registered prerequisite joint
accuracy floors of $0.80$ for $N=2$ and $0.40$ for $N=3$.
M1 achieved only $0.0290$ ($58/2{,}000$) for $N=2$ across all delays, and
$0.0275$--$0.0280$ ($55$--$56/2{,}000$) for $N=3$.
Every cell therefore remains \texttt{FLOOR\_LIMITED\_COMPARISON}.

M2a preserves the dynamic novelty-driven gate and resets $u,r$ to zero
at every active-event offset; it is not the constant-gate arm.
M2b and M2d use the frozen mean gate $A_0=0.061886726$ (displayed rounded),
with M2d also resetting $u,r$. The cell-specific TOST comparison of
M1 and M2a uses equivalence margin $\pm0.005$. The reported $90\%$
confidence intervals lie numerically
inside that margin ($N=2$: $[-0.001515,0.001515]$;
$N=3$: $[-0.002978,0.000841]$ at $0\,\text{ms}$ and
$[-0.002262,0.001186]$ at $5,10\,\text{ms}$).
However, the failed baseline prerequisite prevents a confirmatory
equivalence or functioning-memory-redundancy conclusion in every cell
(\texttt{redundancy\_inferred = False}).

\paragraph{Spiking channel limitation.}
M4\_Count and M4\_Latency use the same physical carrier and query simulation
as M1. Their decoders consume only the spike indicators at cue steps
1--99; the deadline spike at step 100 is excluded, and no analog state
or input-current bypass is provided. Silence is abstention with the
registered zero fallback. The separate composition-risk criterion for
\texttt{DEMONSTRATED\_CHANNEL\_LIMITATION} is
$\mathrm{LCB}_{95\%}(E_{\mathrm{comp}})\ge0.05$, not the operand-delivery
tolerance. Both tested end-to-end readouts meet that criterion across
all cells ($E_{\mathrm{comp}}\in[2.21,2.57]$).
Delay-pooled $95\%$ bootstrap confidence intervals were:
$N=2$: M4\_Count $[2.393533,2.502027]$, M4\_Latency $[2.243582,2.332606]$;
$N=3$: M4\_Count $[2.473506,2.571918]$, M4\_Latency $[2.212398,2.296450]$.
Episode-level abstention rates were $82.1\%$--$83.0\%$
(single-query silence $57.15\%$--$58.60\%$).
These are limitations of the specified carrier-plus-readout implementations;
they do not isolate spike discretization as the sole bottleneck or establish
a universal requirement for population coding.

\end{document}
